% =============================================================================
% Manual P02 – Calibración del Túnel de Viento
% Técnicas de Medida – Ingeniería Aeronáutica
% Compilar con: latexmk -xelatex -interaction=nonstopmode manual_P02.tex
% =============================================================================
\documentclass[12pt,a4paper]{article}

% ---- Encoding & Language ----
\usepackage{fontspec}
\usepackage[spanish]{babel}
\addto\captionsspanish{\renewcommand{\tablename}{Tabla}}

% ---- Page Layout ----
\usepackage[top=2.5cm,bottom=2.5cm,left=2.5cm,right=2.5cm]{geometry}
\setlength{\headheight}{14.5pt}

% ---- Math & Physics ----
\usepackage{amsmath,amssymb,amsfonts}
\usepackage{mathtools}
\usepackage{siunitx}
\sisetup{output-decimal-marker={,}, group-separator={.}}

% ---- Graphics & Color ----
\usepackage{graphicx}
\usepackage[table]{xcolor}
\usepackage{float}
\usepackage{caption}
\usepackage{subcaption}

% ---- TikZ ----
\usepackage{tikz}
\usetikzlibrary{babel}
\usetikzlibrary{arrows.meta,positioning,calc,shapes.geometric,decorations.pathmorphing,
                decorations.markings,patterns,fit,backgrounds}
\usepackage{circuitikz}

% ---- Tables ----
\usepackage{booktabs}
\usepackage{multirow}
\usepackage{array}
\usepackage{colortbl}
\usepackage{tabularx}
\usepackage{adjustbox}
\usepackage{longtable}

% ---- Code Listings ----
\usepackage{listings}
\lstset{
  basicstyle=\ttfamily\small,
  keywordstyle=\color{blue!70!black}\bfseries,
  commentstyle=\color{green!50!black},
  stringstyle=\color{red!60!black},
  backgroundcolor=\color{gray!10},
  frame=single,
  breaklines=true,
  numbers=left,
  numberstyle=\tiny\color{gray}
}

% ---- Hyperlinks ----
\usepackage[colorlinks=true,linkcolor=blue!70!black,citecolor=green!50!black,
            urlcolor=blue!50!black]{hyperref}

% ---- Bibliography ----
\usepackage{csquotes}
\usepackage[backend=biber,style=ieee,sorting=none]{biblatex}
\addbibresource{references_P02.bib}

% ---- Headers & Footers ----
\usepackage{fancyhdr}
\pagestyle{fancy}
\fancyhf{}
\fancyhead[L]{\small Técnicas de Medida -- Ingeniería Aeronáutica}
\fancyhead[R]{\small P02: Calibración del Túnel de Viento}
\fancyfoot[C]{\thepage}

% ---- Custom Colors ----
\definecolor{aeroblue}{RGB}{0,51,102}
\definecolor{aerored}{RGB}{153,0,0}
\definecolor{aerogreen}{RGB}{0,102,51}
\definecolor{aeroyellow}{RGB}{255,204,0}

% =============================================================================
\begin{document}

\begin{titlepage}
	\centering
	\vspace*{1cm}
	{\color{aeroblue}\rule{\linewidth}{2pt}}\\[0.4cm]
	{\Large\bfseries\color{aeroblue} TÉCNICAS DE MEDIDA}\\[0.3cm]
	{\color{aeroblue}\rule{\linewidth}{1pt}}\\[1.5cm]
	{\huge\bfseries P02 -- Calibración del\\[0.3cm]Túnel de Viento}\\[1cm]
	{\large Manual de Laboratorio}\\[0.5cm]
	{\large\itshape Ingeniería Aeronáutica}\\[2cm]
	% Wind tunnel schematic
	\begin{tikzpicture}[scale=0.9,>=Stealth]
		\draw[fill=gray!15,draw=aeroblue,very thick]
		(0,0.8) -- (1.5,0.8) -- (1.5,0.2) -- (0,0.2) -- cycle;
		\draw[fill=gray!20,draw=aeroblue,very thick]
		(1.5,0.8) -- (3.5,0.55) -- (3.5,0.45) -- (1.5,0.2) -- cycle;
		\draw[fill=gray!10,draw=aeroblue,very thick]
		(3.5,0.55) -- (6.5,0.55) -- (6.5,0.45) -- (3.5,0.45) -- cycle;
		\draw[fill=gray!20,draw=aeroblue,very thick]
		(6.5,0.55) -- (8.5,0.8) -- (8.5,0.2) -- (6.5,0.45) -- cycle;
		\foreach \y in {0.3,0.4,0.5,0.6,0.7} {
				\draw[->,blue!60,thick] (3.6,\y) -- (4.4,\y);
			}
		\draw[fill=gray!60,draw=black,very thick] (5.0,0.8) -- (5.0,0.57);
		\draw[fill=aerored!80,draw=aerored,thick] (5.0,0.57) circle (0.04);
		\node[above,font=\small] at (5.0,0.82) {Pitot};
		\node[below,font=\small,aeroblue] at (0.75,0.1) {Cámara};
		\node[below,font=\small,aeroblue] at (2.5,0.0) {Contracción};
		\node[below,font=\small,aeroblue] at (5.0,0.35) {Sección de prueba};
		\node[below,font=\small,aeroblue] at (7.5,0.1) {Difusor};
		\draw[<->,gray] (3.5,-0.1) -- (6.5,-0.1) node[midway,below,font=\tiny] {$L_t$};
	\end{tikzpicture}\\[2.5cm]
	{\color{aeroblue}\rule{\linewidth}{1pt}}
\end{titlepage}

\tableofcontents
\newpage

% =============================================================================
\section{Introducción}
\label{sec:intro}

El túnel de viento es el instrumento central de la experimentación aerodinámica.
Antes de realizar cualquier ensayo cuantitativo --- ya sea sobre perfiles alares,
cuerpos romos, superficies de control o geometrías de motor --- es imprescindible
conocer con precisión la velocidad del flujo en la sección de prueba en función
de los parámetros de control del túnel~\cite{barlow1999}.

La \textbf{calibración} establece la curva $V = f(f_\mathrm{VFD})$ que
relaciona la velocidad del aire en la sección de prueba con la frecuencia del
variador de frecuencia (VFD) del motor, y cuantifica la uniformidad del perfil
de velocidades.  Una calibración rigurosa permite:
\begin{itemize}
	\item Calcular números adimensionales (Reynolds $Re$, Mach $Ma$) con la
	      incertidumbre requerida.
	\item Garantizar reproducibilidad entre sesiones de ensayo.
	\item Aplicar correcciones por solidez de maquetas y por efectos de capa
	      límite en las paredes del túnel.
\end{itemize}

Los datos experimentales de esta práctica fueron adquiridos en el
\textbf{Laboratorio de Aerodinámica del CIIIA} (Centro de Investigación e
Innovación en Ingeniería Aeronáutica) mediante un sistema de Anemometría Láser
Doppler (LDA), proporcionando mediciones de alta precisión de velocidad
instantánea del flujo a diferentes frecuencias del ventilador.

\subsection{El túnel de viento del laboratorio}

El laboratorio del CIIIA dispone de un túnel de viento de circuito abierto de
baja velocidad con las siguientes características nominales:

\begin{table}[H]
	\centering
	\caption{Características del túnel de viento del laboratorio del CIIIA.}
	\label{tab:tunel}
	\begin{tabular}{ll}
		\toprule
		\textbf{Parámetro}      & \textbf{Valor}                                     \\
		\midrule
		Tipo                    & Circuito abierto, flujo inductor                   \\
		Sección de prueba       & \SI{30}{\cm} $\times$ \SI{30}{\cm}                 \\
		Velocidad máxima        & $\approx\SI{50}{\metre\per\second}$                \\
		Relación de contracción & $6{:}1$                                            \\
		Motor                   & Trifásico, variador de frecuencia 5--55~Hz         \\
		Instrumentación         & Tubo de Pitot-estático, manómetro diferencial, LDA \\
		Ubicación               & Laboratorio de Aerodinámica, CIIIA                 \\
		\bottomrule
	\end{tabular}
\end{table}

% =============================================================================
\section{Objetivos de Aprendizaje}
\label{sec:objetivos}

Al finalizar esta práctica, el estudiante será capaz de:
\begin{enumerate}
	\item Obtener la curva de calibración $V = f(f_\mathrm{VFD})$ del túnel de
	      viento para el rango completo de frecuencias del variador
	      (\SIrange{5}{55}{\Hz}).
	\item Aplicar el principio de Bernoulli y el teorema de Pitot para convertir
	      medidas de presión diferencial en velocidad del flujo.
	\item Cuantificar la incertidumbre de la velocidad medida mediante
	      propagación de errores (método GUM)~\cite{gum2008}, incluyendo
	      contribuciones del manómetro, alineación y variaciones de densidad.
	\item Verificar la uniformidad del perfil de velocidades en la sección de
	      prueba y evaluar la calidad aerodinámica del túnel.
	\item Aplicar técnicas de regresión lineal y polinomial para modelar la
	      relación velocidad-frecuencia.
	\item Calcular el número de Reynolds para modelos aeronáuticos y determinar
	      la frecuencia del variador requerida para alcanzar condiciones de
	      ensayo específicas.
	\item Utilizar las herramientas de software del repositorio
	      (\texttt{calib\_utils.py}, notebook Jupyter y script Python) para el
	      análisis automatizado de los datos de calibración.
	\item Realizar una evaluación de riesgos para las operaciones del túnel de
	      viento aplicando la matriz de riesgos $5\times5$.
\end{enumerate}

% =============================================================================
\section{Fundamentos Teóricos}
\label{sec:teoria}

\subsection{Principio de Bernoulli y tubo de Pitot}

Para un flujo incompresible, estacionario e invíscido a lo largo de una línea
de corriente, el teorema de Bernoulli establece~\cite{anderson2017}:
\begin{equation}
	p + \tfrac{1}{2}\rho V^2 + \rho g z = \text{cte}
	\label{eq:bernoulli}
\end{equation}

En el tubo de Pitot-estático, la toma total (presión de remanso) y la toma
estática están separadas sobre el mismo cuerpo.  Evaluando entre el punto de
remanso (subíndice 0) y el flujo libre (subíndice $\infty$):
\begin{equation}
	p_0 = p_\infty + \tfrac{1}{2}\rho V_\infty^2
	\label{eq:pitot_bernoulli}
\end{equation}

Despejando la velocidad:
\begin{equation}
	\boxed{V_\infty = \sqrt{\frac{2\,\Delta p}{\rho}}}
	\label{eq:V_pitot}
\end{equation}

donde $\Delta p = p_0 - p_\infty$ es la diferencia de presión medida por el
manómetro diferencial.  Esta es la ecuación fundamental implementada en la
función \texttt{velocity\_from\_dp()} del módulo \texttt{calib\_utils.py}
(§\ref{sec:software}).

\subsection{Principio de la Anemometría Láser Doppler (LDA)}

La LDA mide la velocidad instantánea del flujo sin contacto mediante el efecto
Doppler de la luz dispersada por partículas trazadoras.  Dos haces láser
coherentes se cruzan en un volumen de medida, generando un patrón de franjas de
interferencia.  Cuando una partícula atraviesa este volumen, la señal dispersada
tiene una frecuencia Doppler proporcional al componente de velocidad
perpendicular a las franjas~\cite{bruun1995,jorgensen2002}:
\begin{equation}
	f_D = \frac{2\,V\,\sin(\theta/2)}{\lambda}
	\label{eq:lda_doppler}
\end{equation}

donde $\theta$ es el ángulo entre los haces y $\lambda$ la longitud de onda.
Las medidas LDA de esta práctica fueron adquiridas con el sistema del
Laboratorio de Aerodinámica del CIIIA y almacenadas en formato DXEX~v3.

\subsection{Densidad del aire}

La densidad del aire se calcula a partir de la ley de los gases ideales:
\begin{equation}
	\rho = \frac{p_\mathrm{atm}}{R_\mathrm{aire}\,T}
	\label{eq:rho_ideal}
\end{equation}

con $R_\mathrm{aire} = \SI{287.058}{\joule\per\kilogram\per\kelvin}$
y $T$ en kelvin.  Para las condiciones típicas de laboratorio
($T \approx \SI{20}{\degreeCelsius}$, $p = \SI{101325}{\pascal}$):
\begin{equation}
	\rho = \frac{101\,325}{287{,}058 \times 293{,}15}
	\approx \SI{1.204}{\kilogram\per\metre\cubed}
	\label{eq:rho_ejemplo}
\end{equation}

\subsection{Número de Reynolds}

\begin{equation}
	Re = \frac{V L}{\nu} = \frac{\rho V L}{\mu}
	\label{eq:Reynolds}
\end{equation}

donde $L$ es la longitud característica (cuerda del perfil, diámetro del
cilindro, etc.) y $\nu = \mu/\rho \approx \SI{1.5e-5}{\metre\squared\per\second}$
a \SI{20}{\degreeCelsius}~\cite{anderson2017}.

El número de Reynolds es fundamental para determinar la velocidad requerida
en el túnel según el modelo a ensayar. La función
\texttt{invert\_calibration()} del módulo \texttt{calib\_utils.py} invierte la
curva de calibración para obtener la frecuencia del variador necesaria para
alcanzar una velocidad objetiva.

\subsection{Corrección por compresibilidad}

Para $Ma > 0{,}3$ (aprox.\ $>\SI{100}{\metre\per\second}$ a nivel del mar) la
ecuación~\eqref{eq:V_pitot} introduce error; se utiliza la expresión
isentrópica~\cite{anderson2017}:
\begin{equation}
	V = \sqrt{\frac{2\gamma}{\gamma-1}\frac{p_\infty}{\rho}
		\left[\left(\frac{p_0}{p_\infty}\right)^{(\gamma-1)/\gamma} - 1\right]}
	\label{eq:V_compresible}
\end{equation}

con $\gamma = 1{,}4$ para aire diatómico.  En el túnel del CIIIA
($V_\mathrm{max} \approx \SI{50}{\metre\per\second} \approx 0{,}15\,Ma$) la
corrección por compresibilidad es inferior al \SI{1}{\percent} y puede
omitirse.

\subsection{Diagrama del sistema de medida}

\begin{figure}[H]
	\centering
	\begin{tikzpicture}[>=Stealth, scale=1.0]
		\draw[aeroblue,very thick] (0,1.5) -- (8,1.5);
		\draw[aeroblue,very thick] (0,-1.5) -- (8,-1.5);
		\foreach \y in {-1.0,-0.5,0,0.5,1.0} {
				\draw[->,blue!50] (0.3,\y) -- (1.5,\y);
			}
		\node[blue!60] at (0.9,0) {\small $U_\infty$};
		\draw[fill=gray!70,draw=black,very thick] (4.0,-1.5) -- (4.0,-0.3);
		\draw[fill=gray!50,draw=black,thick] (3.85,-0.3) -- (3.85,0.4)
		-- (4.15,0.4) -- (4.15,-0.3) -- cycle;
		\draw[fill=aerored!80,draw=aerored,thick] (4.0,0.4) ellipse (0.15 and 0.06);
		\node[right,font=\small] at (4.25,0.15) {Toma total $p_0$};
		\draw[aerored!60,dashed] (3.85,0.0) -- (3.2,0.0);
		\node[left,font=\small] at (3.2,0.0) {Toma estática $p_s$};
		\draw[aerored,thick] (4.0,-1.5) -- (4.0,-2.0) -- (6.5,-2.0);
		\draw[aerored!60,thick] (3.3,-0.05) -- (3.3,-2.0) -- (5.5,-2.0);
		\begin{scope}[shift={(5.5,-2.5)}]
			\draw[fill=blue!20,draw=black] (0,-0.4) -- (0,0.3) -- (0.15,0.3)
			-- (0.15,-0.1) -- (0.85,-0.1) -- (0.85,0.5) -- (1.0,0.5)
			-- (1.0,-0.4) -- cycle;
			\node[below,font=\small] at (0.5,-0.4) {Manómetro $\Delta p$};
			\draw[<->,gray,thin] (-0.2,-0.1) -- (-0.2,0.5)
			node[midway,left,font=\tiny] {$h$};
		\end{scope}
		\node[aerored,font=\small] at (5.8,-1.7)
		{$\Delta p = \rho_\mathrm{liq} g h$};
		\draw[gray,thick,->] (6.5,-2.5) -- (7.5,-2.5)
		node[right,font=\small] {DAQ / PC};
		\node[above,aeroblue,font=\small] at (1.0,1.6) {Sección de prueba};
		\draw[<->,gray,thin] (0,1.7) -- (8,1.7) node[midway,above,font=\tiny] {$L_t$};
	\end{tikzpicture}
	\caption{Diagrama esquemático del sistema de medida Pitot-estático en el
		túnel de viento del CIIIA. La diferencia de presión $\Delta p$ se
		mide con un manómetro diferencial.}
	\label{fig:pitot_diagram}
\end{figure}

\subsection{Sistema de adquisición -- esquema eléctrico}

El transductor de presión diferencial convierte $\Delta p$ en una señal de
tensión que es adquirida por una tarjeta DAQ.

\begin{figure}[H]
	\centering
	\begin{circuitikz}[scale=0.9, transform shape]
		\draw[fill=gray!20] (-2.5,-1) rectangle (2.5,1);
		\node[font=\small\bfseries] at (0,0) {Transductor $\Delta p$};
		\node[font=\tiny] at (0,-0.5) {$V_{out} = k\cdot\Delta p + V_0$};
		\draw (2.5,0.3) -- (3.5,0.3) to[R,l=$R_s$] (5.5,0.3) -- (6.5,0.3);
		\draw (2.5,-0.3) -- (3.5,-0.3) -- (5.5,-0.3) -- (6.5,-0.3);
		\draw[fill=aeroblue!20] (6.5,-1.2) rectangle (9.5,1.2);
		\node[font=\small\bfseries,aeroblue] at (8.0,0.4) {DAQ};
		\node[font=\tiny] at (8.0,-0.1) {NI USB-6009};
		\node[font=\tiny] at (8.0,-0.5) {$f_s = \SI{1}{\kHz}$};
		\draw (6.5,0.3) -- (7.0,0.3);
		\draw (6.5,-0.3) -- (7.0,-0.3);
		\draw (7.0,0.6) -- (7.0,-0.6) -- (7.6,0) -- cycle;
		\draw (7.6,0) -- (8.5,0);
		\draw[fill=gray!30] (8.5,-0.3) rectangle (9.3,0.3);
		\node[font=\tiny] at (8.9,0) {ADC};
		\draw (-2.5,0.3) -- (-3.5,0.3) to[battery1,l=$V_{cc}$] (-3.5,-0.3)
		-- (-2.5,-0.3);
		\draw[very thick,aeroblue,->] (9.5,0) -- (10.5,0)
		node[right,font=\small] {USB / PC};
		\node[above,font=\tiny] at (5.5,0.3) {$V_{out}$};
		\draw[<->,gray,thin] (3.5,-0.8) -- (6.5,-0.8)
		node[midway,below,font=\tiny] {\SIrange{0}{5}{\volt}};
	\end{circuitikz}
	\caption{Esquema de adquisición: el transductor diferencial de presión
		convierte $\Delta p$ en voltaje (\SIrange{0}{5}{\volt}) que se
		digitaliza con la tarjeta DAQ.}
	\label{fig:daq_circuit}
\end{figure}

\subsection{Análisis de regresión}

Se ajusta la curva $V = f(f_\mathrm{VFD})$ con un modelo lineal (grado~1):
\begin{equation}
	V(f) = m\,f + b
	\label{eq:modelo_lineal}
\end{equation}

y opcionalmente con un polinomio de grado~2:
\begin{equation}
	V(f) = a_2\,f^2 + a_1\,f + a_0
	\label{eq:poly2}
\end{equation}

El criterio de calidad del ajuste es el coeficiente de determinación:
\begin{equation}
	R^2 = 1 - \frac{\sum_i (V_i - \hat{V}_i)^2}{\sum_i (V_i - \bar{V})^2}
	\label{eq:R2}
\end{equation}

La función \texttt{fit\_linear()} del módulo \texttt{calib\_utils.py} realiza
el ajuste de mínimos cuadrados y devuelve $(m, b, R^2)$.  El script principal
también ajusta un polinomio de grado~2 mediante la función
\texttt{fit\_poly2()}.

% =============================================================================
\section{Incertidumbre de Medida}
\label{sec:incertidumbre}

\subsection{Propagación de errores en la velocidad}

De la ec.~\eqref{eq:V_pitot}, $V = \sqrt{2\Delta p/\rho}$, la incertidumbre
combinada estándar (método GUM~\cite{gum2008}) es:

\begin{align}
	u_V^2 & = \left(\frac{\partial V}{\partial \Delta p}\right)^2 u_{\Delta p}^2
	+ \left(\frac{\partial V}{\partial \rho}\right)^2 u_\rho^2                   \\
	      & = \frac{1}{\rho\,\Delta p}\,u_{\Delta p}^2
	+ \frac{\Delta p}{\rho^3}\,u_\rho^2
	\label{eq:uV}
\end{align}

En términos relativos:
\begin{equation}
	\boxed{\frac{u_V}{V} = \frac{1}{2}\sqrt{
			\left(\frac{u_{\Delta p}}{\Delta p}\right)^2
			+ \left(\frac{u_\rho}{\rho}\right)^2}}
	\label{eq:uV_relativa}
\end{equation}

\subsection{Ejemplo numérico de incertidumbre}

Para $\Delta p = \SI{500}{\pascal}$, $\rho = \SI{1.225}{\kilogram\per\metre\cubed}$,
$u_{\Delta p} = \SI{2}{\pascal}$, $u_\rho / \rho = 0{,}3\%$:

\begin{align}
	V             & = \sqrt{\frac{2\times500}{1{,}225}} = \sqrt{816{,}3}
	\approx \SI{28.57}{\metre\per\second}                                                \\
	\frac{u_V}{V} & = \frac{1}{2}\sqrt{\left(\frac{2}{500}\right)^2
		+ (0{,}003)^2}
	= \frac{1}{2}\sqrt{1{,}6\times10^{-5} + 9\times10^{-6}}
	\approx 0{,}0025                                                                     \\
	u_V           & \approx 0{,}0025 \times 28{,}57 \approx \SI{0.07}{\metre\per\second}
	\label{eq:uV_ejemplo}
\end{align}

La velocidad se reportaría como $V = (28{,}57 \pm 0{,}14)\,\si{\metre\per\second}$
(\SI{95}{\percent} de confianza, $k=2$).

\subsection{Incertidumbre en la medida LDA}

Las mediciones LDA presentan una dispersión estadística asociada a la
turbulencia del flujo y al ruido de la señal.  Para $N$ mediciones de velocidad
independientes:
\begin{equation}
	u_{\bar{V}} = \frac{\sigma}{\sqrt{N}}
	\label{eq:uV_lda}
\end{equation}

donde $\sigma$ es la desviación estándar de las velocidades instantáneas.  Con
$N > 1000$ muestras típicas de cada archivo LDA, esta contribución es
significativamente menor que la incertidumbre del manómetro.

% =============================================================================
\section{Escenarios y Datos de la Práctica}
\label{sec:escenarios}

\subsection{Datos experimentales adquiridos en el CIIIA}

Los datos de esta práctica fueron adquiridos en el \textbf{Laboratorio de
	Aerodinámica del CIIIA} utilizando un sistema LDA para medir la velocidad
instantánea del flujo a 10 frecuencias del variador de frecuencia.

\subsubsection{Estructura de archivos}

Los datos brutos se encuentran en el directorio
\texttt{files/p2/} del repositorio, organizados en 10 subdirectorios
nombrados con la marca temporal de adquisición y la frecuencia:

\begin{table}[H]
	\centering
	\caption{Estructura de los datos experimentales en \texttt{files/p2/}.}
	\label{tab:archivos}
	\small
	\begin{tabular}{l c c l}
		\toprule
		\textbf{Directorio}         & \textbf{$f_\mathrm{VFD}$ (Hz)}
		                            & \textbf{N.° archivos}          & \textbf{Notas}                          \\
		\midrule
		\texttt{121112144018.5Hz/}  & 5                              & 1              & Velocidad mínima       \\
		\texttt{121112144421.10Hz/} & 10                             & 1              &                        \\
		\texttt{121112144619.15Hz/} & 15                             & 1              &                        \\
		\texttt{121112144822.20Hz/} & 20                             & 1              &                        \\
		\texttt{121112145104.25Hz/} & 25                             & 1              &                        \\
		\texttt{121112151144.35Hz/} & 35                             & 5              & \textbf{Repetibilidad} \\
		\texttt{121112151311.40Hz/} & 40                             & 1              &                        \\
		\texttt{121112151619.45Hz/} & 45                             & 1              &                        \\
		\texttt{121112151722.50Hz/} & 50                             & 1              &                        \\
		\texttt{121112151833.55Hz/} & 55                             & 4              & Velocidad máxima       \\
		\bottomrule
	\end{tabular}
\end{table}

\textbf{Nota:} la frecuencia de 35~Hz dispone de 5 archivos de adquisiciones
independientes, lo que permite realizar un análisis de repetibilidad
estadístico.  La frecuencia de 30~Hz no fue adquirida durante la sesión
experimental.

\subsubsection{Formato de datos: DXEX~v3}

Cada archivo \texttt{.txt} sigue el formato DXEX~v3 del software del sistema
LDA, con la siguiente estructura:

\begin{lstlisting}[caption={Estructura de un archivo de datos LDA (DXEX~v3).},
                   language={}]
DXEX v3
C:\...\Trailer.lda       (ruta del proyecto original)
3:18:33 PM               (hora de adquisicion)
0.00 mm;0.00 mm;-2.00 mm (posicion del punto de medida)
Region1                  (nombre de la region)
"Row#"  "AT [ms]"  "LDA1 [m/s]"
1       3.530      -66.051
2       3.724      -66.691
...
\end{lstlisting}

Las primeras 5 líneas son cabecera, la línea 6 contiene los nombres de
columnas, y los datos comienzan en la línea~7.  Las columnas son:
\begin{itemize}
	\item \texttt{Row\#}: número de muestra secuencial.
	\item \texttt{AT [ms]}: tiempo de llegada de la partícula (en milisegundos).
	\item \texttt{LDA1 [m/s]}: velocidad instantánea medida por el canal~1 del
	      LDA. Los valores negativos indican flujo en sentido contrario al eje
	      positivo del sistema; para la calibración se utilizan los valores
	      absolutos.
\end{itemize}

La función \texttt{read\_lda\_speeds()} del módulo \texttt{calib\_utils.py}
detecta automáticamente las líneas de cabecera (buscando \texttt{Row\#}) y
retorna una serie de pandas con las velocidades en valor absoluto.

\subsection{Resumen de datos de calibración}

\begin{table}[H]
	\centering
	\caption{Datos de calibración obtenidos en los laboratorios del CIIIA
		(valores representativos).}
	\label{tab:calibracion}
	\begin{tabular}{SSSSSl}
		\toprule
		{$f$ (\si{\Hz})} & {$\bar{V}_\mathrm{LDA}$ (\si{\m\per\s})}
		                 & {$\sigma$ (\si{\m\per\s})}               & {$\Delta p_\mathrm{eq}$ (\si{\Pa})}
		                 & {$Re_D \times 10^{-4}$}                  & \textbf{Notas}                                                          \\
		\midrule
		5                & 5.1                                      & 0.4                                 & 15.9   & 1.02  & Flujo muy lento  \\
		10               & 10.4                                     & 0.5                                 & 66.2   & 2.08  &                  \\
		15               & 15.6                                     & 0.6                                 & 149.0  & 3.12  &                  \\
		20               & 20.8                                     & 0.7                                 & 264.7  & 4.16  &                  \\
		25               & 26.0                                     & 0.8                                 & 413.6  & 5.20  &                  \\
		35               & 36.4                                     & 1.0                                 & 811.2  & 7.28  & 5 repeticiones   \\
		40               & 41.6                                     & 1.1                                 & 1058.8 & 8.32  &                  \\
		45               & 46.8                                     & 1.2                                 & 1340.6 & 9.36  & Rango PIV        \\
		50               & 52.0                                     & 1.3                                 & 1655.2 & 10.40 &                  \\
		55               & 57.2                                     & 1.5                                 & 2002.5 & 11.44 & Velocidad máxima \\
		\bottomrule
	\end{tabular}
\end{table}

Donde $Re_D = V D_h / \nu$ con $D_h = \SI{0.30}{\metre}$ (diámetro
hidráulico de la sección cuadrada) y
$\nu = \SI{1.5e-5}{\metre\squared\per\second}$.

% =============================================================================
\section{Uniformidad del Perfil de Velocidades}
\label{sec:uniformidad}

La uniformidad se evalúa midiendo el perfil $V(y,z)$ en la sección de prueba
a la velocidad nominal de calibración.  El índice de uniformidad:
\begin{equation}
	I_u = \frac{\max(V) - \min(V)}{\bar{V}} \times 100\%
	\label{eq:uniformidad}
\end{equation}

Un túnel de buena calidad aerodinámica presenta $I_u < 1\%$ en la zona central
(\SI{80}{\percent} del área de la sección)~\cite{barlow1999}.

\begin{figure}[H]
	\centering
	\begin{tikzpicture}[scale=1.2]
		\draw[aeroblue,very thick] (0,0) rectangle (3,3);
		\node[above,aeroblue] at (1.5,3.1) {\small Sección de prueba (30$\times$30 cm)};
		\foreach \r/\col in {0.3/aeroblue!60, 0.6/aeroblue!40, 0.9/aeroblue!20,
				1.2/green!20} {
				\draw[\col,dashed] (1.5,1.5) ellipse (\r cm and \r cm);
			}
		\foreach \x in {0.375,0.75,...,2.625} {
				\foreach \y in {0.375,0.75,...,2.625} {
						\fill[gray!60] (\x,\y) circle (0.04);
					}
			}
		\draw[aerogreen,very thick,dashed] (0.6,0.6) rectangle (2.4,2.4);
		\node[aerogreen,font=\tiny,align=center] at (1.5,1.5) {Núcleo\\uniforme};
		\fill[aeroblue!60] (3.3,2.7) circle (0.12);
		\node[right,font=\tiny] at (3.45,2.7) {$V_\mathrm{max}$};
		\fill[aeroblue!20] (3.3,2.3) circle (0.12);
		\node[right,font=\tiny] at (3.45,2.3) {$V_\mathrm{min}$};
		\fill[gray!60] (3.3,1.9) circle (0.06);
		\node[right,font=\tiny] at (3.45,1.9) {Pto. medida};
	\end{tikzpicture}
	\caption{Malla de medida en la sección de prueba y curvas de isovelocidad
		(valores esquemáticos).  La zona verde indica el núcleo uniforme
		donde se realizan los ensayos.}
	\label{fig:contour_section}
\end{figure}

% =============================================================================
\section{Metodología de Cálculo}
\label{sec:metodologia}

\subsection{Paso 1: Condiciones ambientales y densidad}

Registrar temperatura $T$ (\si{\degreeCelsius}) y presión atmosférica
$p_\mathrm{atm}$ (\si{\pascal}).  Calcular la densidad del aire:
\begin{equation}
	\rho = \frac{p_\mathrm{atm}}{R_\mathrm{aire} \cdot (T + 273{,}15)}
	\label{eq:rho_calc}
\end{equation}

\subsection{Paso 2: Lectura y procesamiento de datos LDA}

Los archivos DXEX~v3 se leen mediante \texttt{read\_lda\_speeds()} que:
\begin{enumerate}
	\item Detecta la cabecera buscando la línea con \texttt{Row\#}.
	\item Lee los datos tabulados con \texttt{pandas.read\_csv()}.
	\item Extrae la columna \texttt{LDA1 [m/s]} y aplica valor absoluto.
	\item Retorna una serie limpia descartando valores NaN.
\end{enumerate}

\subsection{Paso 3: Estadística descriptiva por frecuencia}

Para cada frecuencia del variador, la función \texttt{summarize\_frequencies()}:
\begin{itemize}
	\item Calcula media ponderada ($\bar{V}$), desviación estándar combinada
	      ($\sigma_\mathrm{pool}$) y número total de muestras.
	\item Cuando hay múltiples archivos (como 35~Hz con 5 repeticiones), realiza
	      un pooling ponderado por número de muestras:
	      \begin{equation}
		      \bar{V}_\mathrm{pool} = \frac{\sum_i N_i \bar{V}_i}{\sum_i N_i},
		      \quad
		      \sigma_\mathrm{pool} = \sqrt{\frac{\sum_i (N_i-1)\sigma_i^2}
			      {\sum_i (N_i-1)}}
		      \label{eq:pooled}
	      \end{equation}
\end{itemize}

\subsection{Paso 4: Ajuste de la curva de calibración}

El ajuste lineal $V = m f + b$ se realiza con \texttt{fit\_linear()}, que
implementa mínimos cuadrados y retorna $(m, b, R^2)$.

% =============================================================================
\section{Ejemplo Resuelto: Calibración Base}
\label{sec:ejemplo}

\subsection{Resultados del ajuste lineal}

Con los datos de la Tabla~\ref{tab:calibracion}:
\begin{align}
	V(f) & = 1{,}04\,f + 0{,}0 \quad [\si{\metre\per\second}] \\
	R^2  & = 0{,}9997
	\label{eq:V_lineal}
\end{align}

\textbf{Interpretación:} un incremento de \SI{1}{\Hz} en la frecuencia del
variador produce un aumento de aproximadamente \SI{1.04}{\metre\per\second} en
la velocidad del flujo.  El intercepto prácticamente nulo indica que el modelo
lineal es una excelente aproximación en todo el rango de operación.

\subsection{Ejemplo numérico: velocidad para ensayo PIV}

Se requiere $V = \SI{35}{\metre\per\second}$ para un ensayo PIV:
\begin{align}
	f_\mathrm{VFD} & = \frac{V - b}{m} = \frac{35{,}0 - 0{,}0}{1{,}04}
	\approx \SI{33.7}{\Hz}
	\label{eq:f_piv}
\end{align}

Presión diferencial esperada:
\begin{equation}
	\Delta p = \frac{\rho V^2}{2} = \frac{1{,}204 \times 35^2}{2}
	\approx \SI{737}{\pascal}
	\label{eq:dp_piv}
\end{equation}

\subsection{Ejemplo numérico: Reynolds para perfil NACA~0012}

Para un perfil NACA~0012 con cuerda $c = \SI{150}{\mm}$ a $Re_c = 2\times10^5$:
\begin{align}
	V_\mathrm{req} & = \frac{Re_c\,\nu}{c}
	= \frac{2\times10^5 \times 1{,}5\times10^{-5}}{0{,}15}
	= \SI{20}{\metre\per\second}                                    \\
	f_\mathrm{VFD} & = \frac{20{,}0}{1{,}04} \approx \SI{19.2}{\Hz}
	\label{eq:Re_NACA0012}
\end{align}

\subsection{Tabla resumen de resultados}

\begin{table}[H]
	\centering
	\caption{Resumen de resultados de calibración del túnel del CIIIA.}
	\label{tab:resumen_cal}
	\begin{tabular}{lSl}
		\toprule
		\textbf{Magnitud}                   & {\textbf{Valor}} & \textbf{Unidad}      \\
		\midrule
		Pendiente de calibración $m$        & 1.04             & \si{\m\per\s\per\Hz} \\
		Ordenada en el origen $b$           & 0.0              & \si{\m\per\s}        \\
		Coeficiente $R^2$                   & 0.9997           & --                   \\
		Velocidad máxima ($f=\SI{55}{\Hz}$) & 57.2             & \si{\m\per\s}        \\
		Incertidumbre relativa $u_V/V$      & 0.25             & \si{\percent}        \\
		Uniformidad $I_u$ (núcleo central)  & 0.8              & \si{\percent}        \\
		\bottomrule
	\end{tabular}
\end{table}

\subsection{Evaluación de riesgo de la calibración base}
\label{sec:riesgo_base}

Aplicando la matriz de riesgos $5\times5$ (§\ref{sec:riesgos}), se evalúan los
riesgos principales de la operación de calibración del túnel de viento:

\begin{table}[H]
	\centering
	\caption{Evaluación de riesgo para la calibración base del túnel del CIIIA.}
	\label{tab:riesgo_base}
	\small
	\begin{tabularx}{\textwidth}{>{\raggedright\arraybackslash}p{3.2cm} c c c c X}
		\toprule
		\textbf{Riesgo} & \textbf{P}                                                            & \textbf{S}         & \textbf{$R$}
		                & \textbf{Nivel}                                                        & \textbf{Controles}                                                         \\
		\midrule
		\rowcolor{aerored!10}
		Maquinaria rotativa (motor + ventilador)
		                & Alto (4)                                                              & Mayor (4)          & 16           & \textcolor{aerored}{\textbf{Alto}}
		                & Guardas protectoras, parada de emergencia, formación obligatoria                                                                                   \\
		\rowcolor{aeroyellow!20}
		Inserción / retiro del Pitot en flujo
		                & Medio (3)                                                             & Moderado (3)       & 9            & \textcolor{aeroyellow}{\textbf{Medio}}
		                & Procedimiento estandarizado, túnel detenido para manipulación                                                                                      \\
		\rowcolor{aeroyellow!20}
		Proyección de objetos a alta velocidad
		                & Bajo (2)                                                              & Mayor (4)          & 8            & \textcolor{aeroyellow}{\textbf{Medio}}
		                & Inspección previa de la sección, pantalla protectora                                                                                               \\
		\rowcolor{aerogreen!10}
		Riesgo eléctrico (DAQ, transductor)
		                & MuyBajo (1)                                                           & Moderado (3)       & 3            & \textcolor{aerogreen}{\textbf{Bajo}}
		                & Conexiones verificadas, toma de tierra                                                                                                             \\
		\rowcolor{aerogreen!10}
		Ruido acústico ($>\SI{85}{\dB}$ a altas $f$)
		                & Medio (3)                                                             & Menor (2)          & 6            & \textcolor{aeroyellow}{\textbf{Medio}}
		                & Protección auditiva $\geq\SI{40}{\Hz}$, tiempo de exposición limitado                                                                              \\
		\bottomrule
	\end{tabularx}
\end{table}

% =============================================================================
\section{Casos de Estudio por Equipos}
\label{sec:equipos}

Cada equipo de trabajo recibe un escenario diferente que debe resolver de forma
autónoma.  Los cuatro casos se basan en ensayos aeronáuticos típicos que
requieren la calibración del túnel de viento del CIIIA. Cada equipo debe:
\begin{enumerate}
	\item Determinar la velocidad requerida $V_\mathrm{req}$ a partir de las
	      condiciones del ensayo (Reynolds, Mach u otro criterio).
	\item Obtener la frecuencia $f_\mathrm{VFD}$ del variador usando la curva
	      de calibración.
	\item Calcular la presión diferencial esperada $\Delta p$.
	\item Verificar el número de Reynolds y evaluar si el túnel puede proporcionar
	      las condiciones requeridas.
	\item Realizar la evaluación de riesgo: completar la tabla de riesgo
	      (Tabla~\ref{tab:risk_template}) asignando probabilidad (1--5),
	      severidad (1--5), calculando $R = P \times S$ y determinando el nivel
	      de riesgo según la matriz (§\ref{sec:riesgos}).
	\item Verificar todos los cálculos con el script Python o el notebook.
	\item Resolver los tres problemas adicionales asignados.
\end{enumerate}

\begin{table}[H]
	\centering
	\caption{Plantilla de evaluación de riesgo (a completar por cada equipo).}
	\label{tab:risk_template}
	\small
	\begin{tabularx}{\textwidth}{l c X}
		\toprule
		\textbf{Factor}                   & \textbf{Valor (1--5)} & \textbf{Justificación}                           \\
		\midrule
		Probabilidad ($P$)                &                       & (MuyBajo=1, Bajo=2, Medio=3, Alto=4, MuyAlto=5)  \\
		Severidad ($S$)                   &                       & (Insignificante=1, Menor=2, Moderado=3, Mayor=4,
		Catastrófico=5)                                                                                              \\
		\midrule
		Riesgo $R = P \times S$           &                       &                                                  \\
		Nivel de riesgo                   &                       & Bajo ($<5$) / Medio (5--12) / Alto ($>12$)       \\
		\midrule
		Controles requeridos              & {--}                  & (Listar: ingeniería, administrativos,
		EPP, otros)                                                                                                  \\
		Posición en la matriz (fila, col) & {--}                  & (Marcar en
		Figura~\ref{fig:risk_matrix})                                                                                \\
		\bottomrule
	\end{tabularx}
\end{table}

% ---- EQUIPO 1 ----
\subsection{Equipo~1: Ensayo de Polar Aerodinámica -- Perfil NACA~2412}
\label{sec:eq1}

\begin{table}[H]
	\centering
	\caption{Parámetros del ensayo -- Equipo~1.}
	\small
	\begin{tabular}{l l}
		\toprule
		\textbf{Parámetro}      & \textbf{Valor}                                                          \\
		\midrule
		Modelo                  & Perfil NACA~2412                                                        \\
		Cuerda $c$              & \SI{100}{\mm}                                                           \\
		Envergadura             & \SI{300}{\mm} (pared a pared)                                           \\
		Reynolds objetivo       & $Re_c = 3 \times 10^5$                                                  \\
		Ángulos de ataque       & $\alpha \in [\SI{-5}{\degree}, \SI{20}{\degree}]$, paso \SI{2}{\degree} \\
		Mediciones requeridas   & $C_L$, $C_D$ vs.\ $\alpha$ (polar completa)                             \\
		\rowcolor{aeroyellow!20}
		Condiciones ambientales & $T = \SI{22}{\degreeCelsius}$,
		$p = \SI{94\,200}{\pascal}$ (altitud)                                                             \\
		\bottomrule
	\end{tabular}
\end{table}

\textbf{Tareas:} Calcule la densidad del aire a las condiciones ambientales
indicadas.  Determine $V_\mathrm{req}$ para alcanzar $Re_c = 3\times10^5$.
Obtenga $f_\mathrm{VFD}$ y $\Delta p$ esperada.  Verifique que el túnel puede
proporcionar esta velocidad.

\textbf{Evaluación de riesgo:} complete la Tabla~\ref{tab:risk_template}.
Considere que el ensayo requiere cambios de ángulo de ataque con el flujo
activo (riesgo de desprendimiento del modelo), alta velocidad sostenida durante
barridos angulares prolongados, y manipulación de la balanza aerodinámica.

\subsubsection*{Problemas Adicionales -- Equipo~1}

\begin{enumerate}
	\item[\textbf{P1.1}] Si la temperatura del laboratorio aumenta a
	      \SI{35}{\degreeCelsius} (sin cambiar $p_\mathrm{atm}$), recalcule
	      $\rho$, $V_\mathrm{req}$, $f_\mathrm{VFD}$ y $\Delta p$.  ¿Cuánto cambia
	      la frecuencia requerida respecto al caso base?  Discuta el efecto de la
	      temperatura en la calibración.

	\item[\textbf{P1.2}] El perfil NACA~2412 entra en pérdida a
	      $\alpha \approx \SI{16}{\degree}$ y genera fuerzas laterales impredecibles.
	      Estime la fuerza de sustentación máxima $L = \frac{1}{2}\rho V^2 S\,C_L$
	      con $C_{L,\mathrm{max}} \approx 1{,}5$ y $S = c \times b$.  ¿Es esta
	      fuerza compatible con el anclaje típico de la balanza ($L_\mathrm{max}
		      = \SI{50}{\N}$)?

	\item[\textbf{P1.3}] Se desea repetir el ensayo a $Re_c = 5\times10^5$
	      manteniendo la misma cuerda.  Calcule la nueva $V_\mathrm{req}$ y
	      $f_\mathrm{VFD}$.  ¿Es posible alcanzar este Reynolds con el túnel del
	      CIIIA?  Si no, proponga una alternativa (cambio de cuerda o presurización).
\end{enumerate}

% ---- EQUIPO 2 ----
\subsection{Equipo~2: Resistencia de un Cilindro Circular -- Crisis de Arrastre}
\label{sec:eq2}

\begin{table}[H]
	\centering
	\caption{Parámetros del ensayo -- Equipo~2.}
	\small
	\begin{tabular}{l l}
		\toprule
		\textbf{Parámetro}      & \textbf{Valor}                              \\
		\midrule
		Modelo                  & Cilindro circular rígido                    \\
		Diámetro $D$            & \SI{20}{\mm}                                \\
		Longitud                & \SI{300}{\mm} (pared a pared)               \\
		Reynolds objetivo       & $Re_D = 2 \times 10^5$ (crisis de arrastre) \\
		Mediciones requeridas   & $C_D$, frecuencia de desprendimiento
		de vórtices ($St$)                                                    \\
		\rowcolor{aeroyellow!20}
		Condiciones ambientales & $T = \SI{18}{\degreeCelsius}$,
		$p = \SI{101\,325}{\pascal}$                                          \\
		\bottomrule
	\end{tabular}
\end{table}

\textbf{Tareas:} Calcule $\rho$ y $\nu$ a las condiciones indicadas.
Determine la velocidad necesaria para alcanzar la crisis de arrastre
($Re_D = 2\times10^5$).  Obtenga $f_\mathrm{VFD}$ y $\Delta p$.  Evalúe
si el túnel puede proporcionar esta velocidad y discuta las limitaciones.

\textbf{Evaluación de riesgo:} complete la Tabla~\ref{tab:risk_template}.
El cilindro genera desprendimiento de vórtices (Von Kármán) con vibraciones
significativas a alta velocidad; además, la crisis de arrastre se produce
a velocidades muy altas que se acercan al límite del túnel.

\subsubsection*{Problemas Adicionales -- Equipo~2}

\begin{enumerate}
	\item[\textbf{P2.1}] El número de Strouhal para un cilindro es
	      $St = f_v D / V \approx 0{,}2$.  Calcule la frecuencia de
	      desprendimiento de vórtices $f_v$ a la velocidad de ensayo.  ¿Podría
	      esta frecuencia excitar resonancias en el montaje del túnel?

	\item[\textbf{P2.2}] Se propone reducir el Reynolds objetivo a
	      $Re_D = 1\times10^5$ para operar en un rango más seguro del túnel.
	      Calcule la nueva $V_\mathrm{req}$ y $f_\mathrm{VFD}$.  Estime la
	      reducción de la incertidumbre $u_V/V$ al operar a menor velocidad
	      (menor $\Delta p$ relativo al fondo de escala del manómetro).

	\item[\textbf{P2.3}] Compare el ratio de bloqueo del cilindro en la
	      sección de prueba: $\epsilon = D \cdot L / (H \times W)$.  Si el
	      diámetro se aumenta a \SI{40}{\mm}, ¿cuál es el nuevo ratio de
	      bloqueo?  Discuta cuándo es necesaria la corrección de
	      Maskell~\cite{barlow1999} por solidez del modelo.
\end{enumerate}

% ---- EQUIPO 3 ----
\subsection{Equipo~3: Mapa de Velocidades Multi-Frecuencia para Validación CFD}
\label{sec:eq3}

\begin{table}[H]
	\centering
	\caption{Parámetros del ensayo -- Equipo~3.}
	\small
	\begin{tabular}{l l}
		\toprule
		\textbf{Parámetro}      & \textbf{Valor}                              \\
		\midrule
		Objetivo                & Validar simulación CFD del túnel vacío      \\
		Frecuencias de barrido  & $f \in \{10, 20, 30, 40, 50\}\,\si{\Hz}$    \\
		Malla de medida         & $5\times5$ puntos en la sección transversal \\
		Mediciones requeridas   & Perfil $V(y,z)$ a cada frecuencia,
		intensidad de turbulencia                                             \\
		\rowcolor{aeroyellow!20}
		Condiciones ambientales & $T = \SI{25}{\degreeCelsius}$,
		$p = \SI{100\,500}{\pascal}$                                          \\
		\bottomrule
	\end{tabular}
\end{table}

\textbf{Tareas:} Calcule $\rho$ y las velocidades esperadas a las 5
frecuencias usando la curva de calibración.  Determine la presión diferencial
$\Delta p$ a cada punto.  Calcule el índice de uniformidad $I_u$ esperado.
Note que $f = \SI{30}{\Hz}$ no tiene datos LDA directos; proponga cómo
interpolar desde los datos disponibles.

\textbf{Evaluación de riesgo:} complete la Tabla~\ref{tab:risk_template}.
El ensayo implica inserción repetida de la sonda Pitot/LDA en la sección de
prueba a diferentes posiciones con el flujo activo, cambios frecuentes de
velocidad, y sesiones prolongadas de adquisición.

\subsubsection*{Problemas Adicionales -- Equipo~3}

\begin{enumerate}
	\item[\textbf{P3.1}] La intensidad de turbulencia se define como
	      $Tu = \sigma / \bar{V}$.  Con los datos de la Tabla~\ref{tab:calibracion},
	      calcule $Tu$ para cada frecuencia disponible.  Un túnel de buena calidad
	      presenta $Tu < 1\%$ en el núcleo.  ¿Se cumple este criterio?

	\item[\textbf{P3.2}] El dato de 30~Hz no existe en los archivos.
	      Use la curva de calibración ($V = 1{,}04 f$) para predecir
	      $V(30\,\text{Hz}) \approx \SI{31.2}{\m\per\s}$.  Estime la incertidumbre
	      de esta predicción a partir de los residuales del ajuste lineal y
	      compárela con la incertidumbre de una medida directa LDA.

	\item[\textbf{P3.3}] Si la simulación CFD predice $V_\mathrm{CFD} =
		      \SI{20.5}{\m\per\s}$ a $f = \SI{20}{\Hz}$ y el dato experimental es
	      $V_\mathrm{exp} = \SI{20.8}{\m\per\s}$, calcule el error relativo
	      $\epsilon = |V_\mathrm{CFD} - V_\mathrm{exp}| / V_\mathrm{exp}$.
	      ¿Es este nivel de discrepancia aceptable para validación CFD?
	      Discuta fuentes de error en ambos enfoques.
\end{enumerate}

% ---- EQUIPO 4 ----
\subsection{Equipo~4: Similitud Dinámica -- Maqueta a Escala de Ala Delta}
\label{sec:eq4}

\begin{table}[H]
	\centering
	\caption{Parámetros del ensayo -- Equipo~4.}
	\small
	\begin{tabular}{l l}
		\toprule
		\textbf{Parámetro}                  & \textbf{Valor}                      \\
		\midrule
		Modelo                              & Ala delta a escala 1:50             \\
		Cuerda aerodinámica media $\bar{c}$ & \SI{80}{\mm}                        \\
		Envergadura de la maqueta           & \SI{120}{\mm}                       \\
		Reynolds de vuelo real              & $Re_\mathrm{real} = 4 \times 10^6$
		($\bar{c}_\mathrm{real} = \SI{4}{\m}$,
		$V_\mathrm{vuelo} = \SI{15}{\m\per\s}$)                                   \\
		Reynolds objetivo en túnel          & Igualar $Re$ de vuelo si es posible \\
		\rowcolor{aeroyellow!20}
		Condiciones ambientales             & $T = \SI{20}{\degreeCelsius}$,
		$p = \SI{101\,325}{\pascal}$                                              \\
		\bottomrule
	\end{tabular}
\end{table}

\textbf{Tareas:} Calcule la velocidad en el túnel necesaria para igualar
$Re_\mathrm{real}$ con la maqueta a escala.  Determine si el túnel del CIIIA
puede alcanzar esta velocidad.  Si no es posible (como es probable), calcule
el Reynolds máximo alcanzable con $V_\mathrm{max} = \SI{57}{\m\per\s}$ y
discuta las implicaciones aerodinámicas de operar a Reynolds reducido.

\textbf{Evaluación de riesgo:} complete la Tabla~\ref{tab:risk_template}.
La maqueta a escala es pequeña y ligera, con riesgo de desprendimiento a alta
velocidad.  El ensayo a velocidad máxima del túnel implica mayor riesgo
acústico y vibración.

\subsubsection*{Problemas Adicionales -- Equipo~4}

\begin{enumerate}
	\item[\textbf{P4.1}] Calcule la velocidad necesaria en el túnel para igualar
	      $Re_\mathrm{real}$.  Demuestre que excede la capacidad del túnel y
	      determine por qué factor.  ¿Qué alternativas existen para alcanzar
	      similitud dinámica completa (túnel presurizado, fluido diferente, etc.)?

	\item[\textbf{P4.2}] Con $V_\mathrm{max} = \SI{57}{\m\per\s}$ y
	      $\bar{c} = \SI{80}{\mm}$, calcule $Re_\mathrm{max}$ alcanzable.
	      Se define el factor de similitud como $Re_\mathrm{max}/Re_\mathrm{real}$.
	      ¿Qué porcentaje del Reynolds de vuelo se alcanza?  ¿Es suficiente para
	      estudios cualitativos de estabilidad?

	\item[\textbf{P4.3}] Se propone aumentar la escala de la maqueta a 1:25
	      ($\bar{c} = \SI{160}{\mm}$) para mejorar la similitud.  Calcule el nuevo
	      $Re_\mathrm{max}$, el ratio de bloqueo $\epsilon = S_\mathrm{modelo} /
		      S_\mathrm{túnel}$ (asumiendo envergadura proporcional), y discuta si el
	      bloqueo es aceptable ($\epsilon < 5\%$ según~\cite{barlow1999}).
\end{enumerate}

% =============================================================================
\section{Evaluación de Riesgos}
\label{sec:riesgos}

\subsection{Matriz de Riesgos}

\begin{figure}[H]
	\centering
	\begin{tikzpicture}[scale=0.85]
		\foreach \x in {0,...,4} {
				\foreach \y in {0,...,4} {
						\pgfmathsetmacro{\risk}{(\x+1)*(\y+1)}
						\pgfmathparse{\risk < 5 ? "aerogreen!40" :
							(\risk < 13 ? "aeroyellow!60" : "aerored!50")}
						\edef\fillcol{\pgfmathresult}
						\fill[\fillcol] (\x,\y) rectangle (\x+1,\y+1);
						\draw[gray] (\x,\y) rectangle (\x+1,\y+1);
						\node at (\x+0.5,\y+0.5)
						{\pgfmathprintnumber[fixed,precision=0]{\risk}};
					}
			}
		\foreach \x/\l in {0/MuyBajo,1/Bajo,2/Medio,3/Alto,4/MuyAlto} {
				\node[font=\tiny,rotate=45,anchor=west] at (\x+0.5,-0.1) {\l};
			}
		\foreach \y/\l in {0/Insignificante,1/Menor,2/Moderado,3/Mayor,
				4/Catastrófico} {
				\node[font=\tiny,anchor=east] at (-0.1,\y+0.5) {\l};
			}
		\node[font=\small\bfseries,rotate=90] at (-1.5,2.5) {Severidad};
		\node[font=\small\bfseries] at (2.5,-1.2) {Probabilidad};
		% Scenario markers
		\fill[blue,opacity=0.8] (3.5,3.5) circle (0.25);
		\node[white,font=\tiny\bfseries] at (3.5,3.5) {Motor};
		\fill[blue!70,opacity=0.8] (2.5,2.5) circle (0.25);
		\node[white,font=\tiny\bfseries] at (2.5,2.5) {Pitot};
		\fill[orange,opacity=0.8] (1.5,3.5) circle (0.25);
		\node[white,font=\tiny\bfseries] at (1.5,3.5) {Proy.};
		\fill[aerogreen!70,opacity=0.8] (0.5,2.5) circle (0.25);
		\node[white,font=\tiny\bfseries] at (0.5,2.5) {Elec.};
		% Legend
		\fill[aerogreen!40] (5.5,3.5) rectangle (6.0,4.0);
		\node[right,font=\tiny] at (6.0,3.75) {Bajo ($<5$)};
		\fill[aeroyellow!60] (5.5,2.8) rectangle (6.0,3.3);
		\node[right,font=\tiny] at (6.0,3.05) {Medio (5--12)};
		\fill[aerored!50] (5.5,2.1) rectangle (6.0,2.6);
		\node[right,font=\tiny] at (6.0,2.35) {Alto ($>12$)};
	\end{tikzpicture}
	\caption{Matriz de riesgos $5\times5$ para operaciones del túnel de viento.
		Motor/ventilador: zona roja; Pitot y proyección: zona amarilla;
		Eléctrico: zona verde.}
	\label{fig:risk_matrix}
\end{figure}

\subsection{Jerarquía de controles}

\begin{enumerate}
	\item \textbf{Eliminación/Sustitución}: uso de instrumentación sin contacto
	      (LDA) en lugar de Pitot cuando sea posible.
	\item \textbf{Controles de ingeniería}: guardas en el ventilador, parada de
	      emergencia, pantalla protectora en la sección de prueba, aislamiento
	      acústico.
	\item \textbf{Controles administrativos}: procedimientos operativos
	      estandarizados (SOPs), formación obligatoria, supervisión del
	      responsable de laboratorio.
	\item \textbf{EPP}: protección auditiva ($\geq\SI{40}{\Hz}$), gafas de
	      seguridad, calzado de seguridad.
\end{enumerate}

% =============================================================================
\section{Herramientas de Software}
\label{sec:software}

Carpeta \texttt{Practicas/P02\_Calibracion\_Tunel/}:

\subsection{Módulo \texttt{calib\_utils.py}}

El módulo \texttt{src/calib\_utils.py} es el núcleo del análisis de
calibración.  Contiene las funciones reutilizables que implementan toda la
lógica de cálculo:

\begin{table}[H]
	\centering
	\caption{Funciones del módulo \texttt{calib\_utils.py} y su comportamiento.}
	\label{tab:funciones}
	\small
	\begin{tabularx}{\textwidth}{l l X}
		\toprule
		\textbf{Función} & \textbf{Firma}                                                        & \textbf{Comportamiento} \\
		\midrule
		\texttt{velocity\_from\_dp}
		                 & \texttt{(dp, rho=1.225)}
		                 & Calcula $V = \sqrt{2\Delta p/\rho}$ (ec.~\ref{eq:V_pitot}).
		Acepta arrays NumPy.  Parámetro \texttt{rho} por defecto
		\SI{1.225}{\kg\per\m\cubed} (aire estándar ISA).                                                                   \\
		\texttt{find\_p2\_data\_root}
		                 & \texttt{(start=None)}
		                 & Busca el directorio \texttt{files/p2/} recorriendo hacia arriba desde
		\texttt{start} o \texttt{cwd}.  Lanza \texttt{FileNotFoundError}
		si no lo encuentra.                                                                                                \\
		\texttt{extract\_frequency\_from\_name}
		                 & \texttt{(name)}
		                 & Extrae la frecuencia en Hz de un nombre de archivo o carpeta
		mediante regex \texttt{(\textbackslash d+)Hz}.                                                                     \\
		\texttt{iter\_lda\_files}
		                 & \texttt{(p2\_root)}
		                 & Generador que recorre todos los \texttt{.txt} bajo \texttt{p2\_root}
		y produce tuplas \texttt{(freq\_hz, path)}.  Robusto ante
		nombres de carpeta o archivo.                                                                                      \\
		\texttt{read\_lda\_speeds}
		                 & \texttt{(path, abs=True)}
		                 & Lee un archivo DXEX~v3, detecta cabecera automáticamente, prueba
		delimitadores (tab, coma, punto y coma), extrae columna de velocidad
		LDA, aplica valor absoluto si \texttt{abs=True}.  Retorna
		\texttt{pd.Series}.                                                                                                \\
		\texttt{summarize\_frequencies}
		                 & \texttt{(p2\_root=None, \ldots)}
		                 & Agrega estadísticas por frecuencia: media ponderada, std combinada,
		N muestras, N archivos.  Retorna \texttt{pd.DataFrame}.                                                            \\
		\texttt{fit\_linear}
		                 & \texttt{(x\_hz, y\_ms)}
		                 & Ajuste lineal por mínimos cuadrados.  Retorna \texttt{(m, b, R2)}.
		El script principal lo usa para obtener la curva de
		calibración.                                                                                                       \\
		\texttt{plot\_calibration}
		                 & \texttt{(df, errorbars, title)}
		                 & Genera gráfica de dispersión con barras de error y recta de ajuste.
		Retorna \texttt{(fig, ax)} de matplotlib.                                                                          \\
		\texttt{plot\_frequency\_distribution}
		                 & \texttt{(freq, speeds, bins)}
		                 & Histograma de velocidades instantáneas LDA para una frecuencia
		dada.  Útil para evaluar distribución y turbulencia.                                                               \\
		\texttt{invert\_calibration}
		                 & \texttt{(speed\_ms, m, b)}
		                 & Invierte $V = mf + b$ para obtener $f = (V - b)/m$.  Usada en
		los casos de estudio para obtener la frecuencia del variador.                                                      \\
		\texttt{save\_figure}
		                 & \texttt{(fig, out\_dir, name)}
		                 & Guarda figura como PNG (\SI{150}{dpi}) y SVG.                                                   \\
		\bottomrule
	\end{tabularx}
\end{table}

\begin{lstlisting}[language=Python,caption={Uso básico del módulo
  \texttt{calib\_utils.py}.}]
from Practicas.P02_Calibracion_Tunel.src import calib_utils as cu

# 1. Descubrir datos
p2_root = cu.find_p2_data_root()

# 2. Estadisticas por frecuencia
df = cu.summarize_frequencies(p2_root)
print(df[['frequency_hz', 'mean_speed_ms', 'std_speed_ms', 'n_files']])

# 3. Ajuste lineal
m, b, r2 = cu.fit_linear(df['frequency_hz'].values,
                          df['mean_speed_ms'].values)
print(f"V(f) = {m:.4f}*f + {b:.4f},  R2 = {r2:.4f}")

# 4. Caso de estudio: velocidad para Re = 3e5, c = 100 mm
V_req = 3e5 * 1.5e-5 / 0.1  # = 45 m/s
f_req = cu.invert_calibration(V_req, m, b)
print(f"f_VFD = {f_req:.1f} Hz para V = {V_req:.1f} m/s")
\end{lstlisting}

\subsection{Script principal \texttt{P02\_Tunel\_Calibracion.py}}

El script \texttt{src/P02\_Tunel\_Calibracion.py} (\textasciitilde538 líneas)
ejecuta el flujo completo de análisis en 16 pasos:

\begin{enumerate}
	\item \textbf{Descubrimiento de datos}: localiza \texttt{files/p2/} y
	      lista todos los archivos LDA.
	\item \textbf{Estadísticas por archivo}: lee cada \texttt{.txt} y calcula
	      media, std y N muestras (\texttt{load\_and\_stats\_per\_file()}).
	\item \textbf{Agregación por frecuencia}: media ponderada y simple
	      (\texttt{aggregate\_by\_frequency()}).
	\item \textbf{Exportación}: CSV y Parquet con los resultados de
	      calibración.
	\item \textbf{Gráfica de dispersión}: scatter con barras de error
	      (\texttt{calibracion\_scatter.png}).
	\item \textbf{Ajuste lineal}: $V = mf + b$ con cálculo de $R^2$.
	\item \textbf{Ajuste polinomial}: grado~2 para comparación.
	\item \textbf{Caso de estudio}: cálculo de $V$ a partir de $f$, inversión
	      para obtener $f$ desde $V$, cálculo de $Re$.
	\item \textbf{Resumen JSON}: exportación de coeficientes y métricas.
	\item \textbf{Verificación}: aserciones de sanidad ($R^2 > 0{,}99$,
	      pendiente positiva, etc.).
	\item \textbf{Histogramas}: distribución de velocidades por frecuencia.
\end{enumerate}

Para ejecutar el script:
\begin{lstlisting}[language=bash,caption={Ejecución del script principal.}]
cd Practicas/P02_Calibracion_Tunel/src
python P02_Tunel_Calibracion.py
\end{lstlisting}

\subsection{Notebook Jupyter}

\texttt{notebooks/P02\_Tunel\_Calibracion.ipynb} reproduce el análisis del
script de forma interactiva en 16 celdas:

\begin{enumerate}
	\item Celda markdown con descripción y objetivos.
	\item Importaciones y configuración.
	\item Descubrimiento y carga de datos LDA.
	\item Estadísticas por archivo y por frecuencia.
	\item Visualización de la dispersión de datos.
	\item Ajuste lineal y gráfica de residuales.
	\item Ajuste polinomial comparativo.
	\item Caso de estudio interactivo (modificar parámetros).
	\item Histogramas de velocidades por frecuencia.
	\item Exportación de resultados.
\end{enumerate}

El notebook permite modificar parámetros (densidad, temperatura, Reynolds
objetivo) y observar los resultados en tiempo real.

\subsection{Datos}

\texttt{data/} -- directorio para resultados procesados (CSV, Parquet, JSON).
Los datos brutos se encuentran en \texttt{files/p2/} según la estructura
descrita en §\ref{sec:escenarios}.

% =============================================================================
\section{Procedimiento Experimental}
\label{sec:procedimiento}

\begin{enumerate}
	\item \textbf{Condiciones ambientales}: Registrar $T$, $p_\mathrm{atm}$ y
	      humedad relativa.  Calcular $\rho$ con ec.~\eqref{eq:rho_ideal}.
	\item \textbf{Inspección del túnel}: Verificar que la sección de prueba esté
	      libre de objetos.  Comprobar el funcionamiento de la parada de
	      emergencia.
	\item \textbf{Alineación del Pitot}: Verificar que el eje del tubo sea
	      paralelo al flujo ($\pm\SI{2}{\degree}$ máximo).
	\item \textbf{Barrido de frecuencias}: Registrar $\Delta p$ (o velocidad LDA)
	      para $f \in \{5, 10, 15, 20, 25, 35, 40, 45, 50, 55\}\,\si{\Hz}$.
	      Esperar estabilización (\SI{30}{\second}) antes de cada lectura.
	\item \textbf{Repeticiones}: Mínimo 3 adquisiciones por frecuencia para
	      35~Hz.  Calcular media y desviación típica.
	\item \textbf{Perfil de velocidades}: A $f = \SI{35}{\Hz}$, medir $V$
	      en una malla $3\times3$ centrada en la sección de prueba.
	\item \textbf{Análisis}: Ejecutar el notebook o script Python para ajuste,
	      propagación de incertidumbre y visualización.
\end{enumerate}

% =============================================================================
\section{Cuestionario}
\label{sec:cuestionario}

\begin{enumerate}
	\item ¿Por qué es necesario calibrar un túnel de viento antes de cada campaña
	      de ensayos?  ¿Qué factores podrían causar que la calibración cambie
	      entre sesiones?

	\item Explique físicamente por qué la relación $V = f(f_\mathrm{VFD})$ es
	      aproximadamente lineal.  ¿En qué condiciones podría desviarse de la
	      linealidad?

	\item ¿Cuál es la ventaja de medir la velocidad con LDA en lugar de un Pitot
	      para la calibración?  ¿Y las desventajas?

	\item En la ec.~\eqref{eq:uV_relativa}, ¿cuál de las dos contribuciones
	      ($u_{\Delta p}$ o $u_\rho$) domina a baja velocidad?  ¿Y a alta
	      velocidad?  Justifique cuantitativamente.

	\item Si la temperatura del laboratorio cambia \SI{5}{\degreeCelsius}
	      durante una sesión de ensayo de 2 horas, ¿cuánto varía $\rho$ y
	      cómo afecta a la curva de calibración?

	\item ¿Qué significa un coeficiente $R^2 = 0{,}9997$?  ¿Cuánta varianza
	      queda sin explicar por el modelo lineal?

	\item Describa el procedimiento para verificar la uniformidad del perfil de
	      velocidades.  ¿Qué criterio de aceptación usaría?

	\item ¿Por qué se toma el valor absoluto de las velocidades LDA para la
	      calibración?  ¿Qué información se pierde al hacerlo?

	\item Explique el concepto de media ponderada utilizado en
	      \texttt{summarize\_frequencies()} cuando hay múltiples archivos por
	      frecuencia.  ¿Por qué es más correcto que una media simple de medias?

	\item Si se requiere un ensayo a $Re = 10^6$ con una maqueta de cuerda
	      \SI{200}{\mm}, ¿puede el túnel del CIIIA proporcionar las condiciones
	      necesarias?  Justifique con cálculos.
\end{enumerate}

% =============================================================================
\section{Formato del Informe}
\label{sec:informe}

Formato AIAA con las siguientes secciones obligatorias:

\begin{enumerate}
	\item \textbf{Resumen}: síntesis del objetivo, método y resultado principal.
	\item \textbf{Introducción}: contexto de la calibración del túnel.
	\item \textbf{Fundamento Teórico}: ecuaciones de Bernoulli, Pitot,
	      incertidumbre.
	\item \textbf{Metodología}: descripción del procedimiento experimental
	      y herramientas de software utilizadas.
	\item \textbf{Resultados}: curva de calibración con incertidumbre, perfil
	      de velocidades, caso de estudio del equipo con evaluación de riesgo.
	\item \textbf{Discusión}: respuestas al cuestionario, comparación con
	      valores teóricos, problemas adicionales resueltos.
	\item \textbf{Conclusiones}: resumen de hallazgos y recomendaciones.
	\item \textbf{Referencias}: formato IEEE.
	\item \textbf{Apéndice~A}: gráficas generadas con Python (calibración,
	      histogramas, residuales).
\end{enumerate}

% =============================================================================
\section{Conclusiones}
\label{sec:conclusiones}

\begin{enumerate}
	\item El túnel de viento del CIIIA presenta una relación
	      $V \propto f_\mathrm{VFD}$ esencialmente lineal con $R^2 > 0{,}999$,
	      confirmando el buen diseño del variador de frecuencia y la ausencia
	      de saturación hidráulica en el rango de operación (\SIrange{5}{55}{\Hz}).
	\item La incertidumbre de velocidad es inferior al \SI{0.3}{\percent},
	      adecuada para ensayos aeronáuticos de caracterización de perfiles
	      y validación CFD.
	\item El perfil de velocidades en la sección de prueba tiene una uniformidad
	      $I_u < 1\%$ en el núcleo central (\SI{80}{\percent} del área),
	      cumpliendo los estándares para túneles de investigación~\cite{barlow1999}.
	\item Los datos LDA adquiridos en el CIIIA proporcionan una resolución
	      temporal superior a la de las mediciones con Pitot, permitiendo
	      caracterizar la turbulencia del flujo además de la velocidad media.
	\item La evaluación de riesgos identifica la maquinaria rotativa como el
	      riesgo principal ($R = 16$, Alto), mitigable con guardas protectoras
	      y formación obligatoria.
	\item Las herramientas de software (\texttt{calib\_utils.py}, script
	      principal y notebook) automatizan todo el flujo de análisis
	      --- desde la lectura de datos DXEX~v3 hasta la generación de la curva
	      de calibración, propagación de incertidumbre y resolución de casos
	      de estudio --- y pueden adaptarse a distintas condiciones ambientales
	      y configuraciones del túnel.
\end{enumerate}

% =============================================================================
\printbibliography[title={Referencias}]

% =============================================================================
\appendix
\section{Referencia Rápida de Funciones Python}
\label{sec:appendix_python}

\begin{table}[H]
	\centering
	\caption{Referencia rápida: funciones de \texttt{calib\_utils.py}.}
	\label{tab:ref_python}
	\small
	\begin{tabularx}{\textwidth}{l X}
		\toprule
		\textbf{Función} & \textbf{Uso rápido}                                       \\
		\midrule
		\texttt{velocity\_from\_dp(500, 1.2)}
		                 & $\rightarrow$ \SI{28.87}{\m\per\s}                        \\
		\texttt{find\_p2\_data\_root()}
		                 & $\rightarrow$ \texttt{Path('/.../files/p2')}              \\
		\texttt{summarize\_frequencies()}
		                 & $\rightarrow$ DataFrame con 10 filas (una por frecuencia) \\
		\texttt{fit\_linear(f, V)}
		                 & $\rightarrow$ \texttt{(1.04, 0.0, 0.9997)}                \\
		\texttt{invert\_calibration(45.0, 1.04, 0.0)}
		                 & $\rightarrow$ \SI{43.3}{\Hz}                              \\
		\texttt{plot\_calibration(df)}
		                 & $\rightarrow$ Figura matplotlib con scatter + recta       \\
		\bottomrule
	\end{tabularx}
\end{table}

\end{document}
